% SPDX-FileCopyrightText: Copyright © 2026 Idiap Research Institute <contact@idiap.ch>
% SPDX-FileContributor: Andre Freitas <andre.freitas@idiap.ch>
% SPDX-License-Identifier: MIT
% Neuro-symbolic AI Group
% File Documentation
%
% Motivation:
%   Protect Mathformer behavior while the framework evolves across Lean, ATP, KB, SymPy, docs, and skills.
%
% Intent:
%   Encode regression checks that catch parser breakage, semantic drift, and missing workflow evidence.
%
% Function:
%   Contains executable tests, fixtures, and manifests used for local and CI validation.
%
% Maintenance Notes:
%   Keep this header synchronized with major responsibility changes, and prefer updating tests or examples when behavior changes.

\documentclass[11pt]{article}
\usepackage[utf8]{inputenc}
\usepackage[T1]{fontenc}
\usepackage{lmodern}
\usepackage{amsmath, amssymb, amsthm}
\usepackage{hyperref}
\usepackage{geometry}
\geometry{margin=1in}

\title{Choquet Theory}
\author{}
\date{}

\newtheorem{theorem}{Theorem}

\begin{document}
\maketitle

\section*{Overview}

Choquet theory is an area of functional analysis and convex analysis concerned with measures supported on the extreme points of a convex set. Roughly speaking, every vector in a suitable convex set should be representable as a weighted average of extreme points, where the notion of weighted average is generalized from a finite convex combination to an integral over the set of extreme points.

Let $C$ be a subset of a real vector space $V$, and let $E$ denote the set of extreme points of $C$. In the infinite-dimensional setting, one studies cases where $V$ is a locally convex Hausdorff topological vector space, extending ideas that are familiar from finite-dimensional convexity.

Choquet theory has also become a general paradigm for studying convex cones via their extreme rays and, more broadly, for analyzing positivity structures in mathematics.

\section*{Finite-dimensional motivation}

For points $v,w \in V$, the line segment joining them is
\[
\{ \lambda v + (1-\lambda) w : 0 \le \lambda \le 1 \}.
\]

A classical result of Minkowski says that in Euclidean space, a bounded closed convex set $C$ is the convex hull of its set $E$ of extreme points. Thus every $c \in C$ can be written as a finite convex combination of points of $E$.

Equivalently, assigning nonnegative weights $w(e)$ to points $e \in E$, with all but finitely many equal to zero, one may write
\[
c = \sum_{e \in E} w(e)\, e,
\qquad
\sum_{e \in E} w(e) = 1.
\]
The coefficients $w(e)$ determine a probability measure supported on a finite subset of $E$. For any affine function $f$ on $C$, one then has
\[
f(c) = \int f(e)\, dw(e).
\]

\section*{Choquet's theorem}

\begin{theorem}[Choquet]
Let $C$ be a compact convex metrizable subset of a locally convex topological vector space $V$. For every $c \in C$, there exists a probability measure $w$ on $C$, supported on the set $E$ of extreme points of $C$, such that for every affine function $f$ on $C$,
\[
f(c) = \int f(e)\, dw(e).
\]
\end{theorem}

In many applications, $V$ may be a Banach space, but the theorem is useful more generally, including settings in which $V$ itself is not metrizable, for example when it is equipped with the weak-$^\ast$ topology.

The Krein--Milman theorem follows from Choquet's theorem. Another corollary is the Riesz representation theorem for states on the continuous functions on a metrizable compact Hausdorff space.

\section*{Nonmetrizable case}

When $C$ is not metrizable, the Choquet--Bishop--de Leeuw theorem provides a related representation: one obtains a probability measure that vanishes on Baire subsets of $C$ containing no extreme points.

\section*{Uniqueness and Choquet simplices}

Besides existence of representing measures, one may ask whether such a measure is unique.

Uniqueness fails in general, even in finite dimensions; examples include a cube or a ball in $\mathbb{R}^3$. However, uniqueness does hold for finite-dimensional simplices. This leads to the notion of a \emph{Choquet simplex}: a convex set in which every point is represented by a unique probability measure on the extreme points.

\section*{See also}

\begin{itemize}
    \item Carath\'eodory's theorem
    \item Helly's theorem
    \item Krein--Milman theorem
    \item Shapley--Folkman lemma
    \item Convex analysis
    \item Functional analysis
\end{itemize}

\section*{References}

\begin{itemize}
    \item Robert R. Phelps, \emph{Lectures on Choquet's Theorem}, Lecture Notes in Mathematics, vol.~1757, Springer, 2001.
    \item Yassine El Maazouz, \emph{On Choquet's Theorem}.
    \item Errett Bishop and Karl de Leeuw, ``The representations of linear functionals by measures on sets of extreme points,'' \emph{Annales de l'Institut Fourier}, 9 (1959), 305--331.
\end{itemize}

\bigskip
\noindent\textbf{Source note.} This \LaTeX{} file is a clean article-style conversion based on the Wikipedia page ``Choquet theory.''

\end{document}
